\section{Installation and Deployment}
\label{app:install}

\subsection{Dependencies}

The system has three parts: a Laravel API (\texttt{backend}) and two React web apps (\texttt{frontend} and \texttt{leodeguita}).
Table~\ref{tab:dependencies} lists what is needed to install and run it.

\begin{table}[H]
  \centering
  \begin{tabularx}{\linewidth}{L{3.3cm}L{3.9cm}Y}
    \toprule
    Component & Version & Purpose \\
    \midrule
    \multicolumn{3}{@{}l}{\textit{Tools on the computer}} \\
    Git & Any recent version & Clone the repository. \\
    PHP & 8.2 or newer (CI uses 8.3) & Run the backend. \\
    PHP extensions & \texttt{mbstring}, \texttt{gd}, \texttt{zip}, \texttt{bcmath}, \texttt{pdo\_sqlite}, \texttt{sqlite3}; \texttt{pdo\_pgsql} for PostgreSQL & The CI installs the first six. The Docker image installs \texttt{pdo\_pgsql} and \texttt{zip}. \\
    Composer & 2 & Install backend libraries. \\
    Node.js and npm & 20 (CI), 26 (Vercel) & Install and run the web apps. \\
    PostgreSQL & 18 on Railway & Main database. Local development can use PostgreSQL or SQLite. Tests use an in-memory SQLite database. \\
    Docker (optional) & Any recent version & Build the backend image from \texttt{backend/Dockerfile}, as Railway does. \\
    \midrule
    \multicolumn{3}{@{}l}{\textit{Main libraries (installed by Composer and npm)}} \\
    Laravel & 12 & Backend framework. \\
    Laravel Sanctum & 4.2 & Token authentication. \\
    Laravel DomPDF & 3.1 & PDF generation. \\
    React, Vite, TypeScript & 19, 7, 5.9 & Web apps. \\
    Vitest and Cypress & 4 and 15 & Tests of the web apps. Cypress is only in \texttt{frontend}. \\
    \midrule
    \multicolumn{3}{@{}l}{\textit{Online accounts (only for deployment)}} \\
    GitHub & Free account & Repository and CI\@. \\
    Railway & Hobby plan & Hosts the API, the PostgreSQL database and a volume for uploaded files. \\
    Vercel & Hobby plan & Hosts the two web apps. \\
    \bottomrule
  \end{tabularx}
  \caption{Dependencies.}
  \label{tab:dependencies}
\end{table}

\subsection{Installation Steps}

These steps run the system on one computer.
They use Linux or macOS commands; on Windows, use \texttt{copy} instead of \texttt{cp}.

\begin{enumerate}[nosep]
  \item \textbf{Get the code.}
        The tags \texttt{sprint1-deploy} and \texttt{sprint2-deploy} mark the delivered version of each sprint.
\begin{lstlisting}[language={}]
git clone https://github.com/Darwin4050E/Leodega.git
cd Leodega
git checkout sprint2-deploy
\end{lstlisting}
  \item \textbf{Install and start the backend.}
        By default the backend uses SQLite, so the empty database file must exist before \texttt{migrate}.
        To use PostgreSQL instead, skip the \texttt{touch} line and set the database variables in \texttt{backend/.env} (Table~\ref{tab:env-vars}).
        The seeder creates the admin user and, outside production, demo data.
\begin{lstlisting}[language={}]
cd backend
composer install
cp .env.example .env
php artisan key:generate
touch database/database.sqlite
php artisan migrate
php artisan db:seed
php artisan storage:link
php artisan serve        # http://127.0.0.1:8000
\end{lstlisting}
  \item \textbf{Install and start the main web app} in a second terminal.
        Create \texttt{frontend/.env} with the line \texttt{VITE\_API\_URL=http://127.0.0.1:8000/api}.
        If the Cypress download fails, run \texttt{npm ci} as \texttt{CYPRESS\_INSTALL\_BINARY=0 npm ci}.
\begin{lstlisting}[language={}]
cd frontend
npm ci
npm run dev              # http://localhost:5173
\end{lstlisting}
  \item \textbf{Install and start the mobile web app} (\texttt{leodeguita}) in a third terminal.
        Create \texttt{leodeguita/.env} with the same \texttt{VITE\_API\_URL} line.
\begin{lstlisting}[language={}]
cd leodeguita
npm ci
npm run dev              # http://localhost:5174
\end{lstlisting}
  \item \textbf{Run the tests} (optional).
        Run them one at a time.
        If the computer is busy, some tests of \texttt{leodeguita} can fail by timeout, because each test has a limit of 5 seconds.
        Expected result for \texttt{sprint2-deploy}: backend 1820 passed and 17 skipped, \texttt{frontend} 1553 passed, and \texttt{leodeguita} 478 passed.
\begin{lstlisting}[language={}]
cd backend && php artisan test
cd frontend && npm test
cd leodeguita && npm test
\end{lstlisting}
\end{enumerate}

\subsection{Environment Variables}

Table~\ref{tab:env-vars} lists the variables the system uses.
The backend reads them from \texttt{backend/.env} locally and from the service variables on Railway.
The web apps read \texttt{VITE\_API\_URL} when they are built.
The example values come from the team's deployment.
Anything that depends on the environment is shown as a placeholder in angle brackets, and real secrets are never written in the repository.

\begin{longtable}{@{}L{4.7cm}L{5.0cm}L{5.1cm}@{}}
  \toprule
  Variable & Purpose & Example value \\
  \midrule
  \endhead
  \multicolumn{3}{@{}l}{\textit{Backend (Laravel)}} \\
  \texttt{APP\_NAME} & Application name. & \texttt{Leodega} \\
  \texttt{APP\_ENV} & Environment. In \texttt{production}, demo data is not created. & \texttt{production} (locally: \texttt{local}) \\
  \texttt{APP\_KEY} & Encryption key. Create it with \texttt{php artisan key:generate}, or add \texttt{-{}-show} to print it for Railway. Keep it secret. & \texttt{base64:<generated>} \\
  \texttt{APP\_DEBUG} & Show detailed errors. Keep it \texttt{false} in production. & \texttt{false} \\
  \texttt{APP\_URL} & Public URL of the API\@. & \texttt{https://}\allowbreak\texttt{<railway-domain>} \\
  \texttt{FRONTEND\_URL} & URL of the main web app. It is used in links of invitations and emails. Default: \texttt{http://localhost:5173}. & \texttt{https://}\allowbreak\texttt{<frontend-domain>} \\
  \texttt{DB\_CONNECTION} & Database driver. Default: \texttt{sqlite}. & \texttt{pgsql} \\
  \texttt{DB\_URL} & Full database address. Used on Railway. & \texttt{\$\{\{Postgres.}\allowbreak\texttt{DATABASE\_URL\}\}} (reference to the Railway database) \\
  \texttt{DB\_HOST}, \texttt{DB\_PORT}, \texttt{DB\_DATABASE}, \texttt{DB\_USERNAME}, \texttt{DB\_PASSWORD} & Connection to a local PostgreSQL, instead of \texttt{DB\_URL}. Default host and port: \texttt{127.0.0.1} and \texttt{5432}. & \texttt{127.0.0.1}, \texttt{5432}, \texttt{leodega}, \texttt{<user>}, \texttt{<password>} \\
  \texttt{TRUSTED\_HOSTS} & Extra host names the API accepts. Comma separated, without \texttt{https://}. & \texttt{<railway-domain>} \\
  \texttt{CORS\_ALLOWED\_ORIGINS} & Web apps allowed to call the API\@. Comma separated, without a final slash. & \texttt{https://}\allowbreak\texttt{<frontend>.vercel.app,} \texttt{https://}\allowbreak\texttt{<leodeguita>.vercel.app} \\
  \texttt{SESSION\_DRIVER}, \texttt{CACHE\_STORE}, \texttt{QUEUE\_CONNECTION} & Where sessions, cache and queued jobs are stored. & \texttt{database} \\
  \texttt{MAIL\_MAILER} & Mail delivery. The value \texttt{log} writes emails in the log instead of sending them. & \texttt{log} \\
  \texttt{LOG\_CHANNEL} & Where the logs go. & \texttt{stderr} \\
  \texttt{PHP\_CLI\_SERVER\_WORKERS} & Optional. Number of workers of the PHP server. Saves memory. & \texttt{2} \\
  \texttt{PORT} & Port of the server. Railway sets it. The Dockerfile uses 8080 when it is missing. & \texttt{8080} \\
  \texttt{PAYMENT\_HOLD\_MINUTES} & Optional. Minutes that a pending reservation holds its dates. Default: 15. & \texttt{15} \\
  \midrule
  \multicolumn{3}{@{}l}{\textit{Web apps (\texttt{frontend} and \texttt{leodeguita})}} \\
  \texttt{VITE\_API\_URL} & Base URL of the API, ending in \texttt{/api}. It is built into the app, so changing it needs a new build. & \texttt{http://127.0.0.1:8000/api} (locally); \texttt{https://}\allowbreak\texttt{<railway-domain>/api} (Vercel) \\
  \bottomrule
  \caption{Environment variables (no secrets).}
  \label{tab:env-vars}
\end{longtable}

\subsection{Deployment}

The API and the database run on Railway.
The two web apps run on Vercel.
The branches \texttt{deploy/sprint1} and \texttt{deploy/sprint2} hold the code of each sprint plus one small deployment commit, and the tags \texttt{sprint1-deploy} and \texttt{sprint2-deploy} mark them.
Deploy Sprint 1 first and then Sprint 2.
Do not go back from Sprint 2 to Sprint 1 on the same database, because Sprint 2 adds database migrations.

\paragraph{Backend and database on Railway.}
\begin{enumerate}[nosep]
  \item Create a project and add a PostgreSQL database service (the Railway template).
  \item Add a service from the GitHub repository.
        Set \textit{Branch} to \texttt{deploy/sprint2} and \textit{Root Directory} to \texttt{/backend}.
        Railway builds the \texttt{Dockerfile} in that folder.
  \item Set the \textit{Pre-deploy Command} to \texttt{php artisan migrate -{}-force}, so each deployment updates the database.
  \item Set the \textit{Healthcheck Path} to \texttt{/up}, and generate a public domain for the service (it listens on port 8080).
  \item Add a volume mounted at \texttt{/app/storage/app}.
        Without it, uploaded photos and permits are lost on every deployment.
  \item Set the environment variables (Table~\ref{tab:env-vars}).
  \item First deployment only: create the admin user.
        Set the pre-deploy command to the line below, deploy once, and then set it back to \texttt{php artisan migrate -{}-force}.
        Demo data is not created in production.
\begin{lstlisting}[language={}]
sh -c "php artisan migrate --force && php artisan db:seed --force"
\end{lstlisting}
\end{enumerate}

\paragraph{Web apps on Vercel.}
Create two projects from the same repository, one with \textit{Root Directory} \texttt{frontend} and one with \texttt{leodeguita}.
In both, use the Vite preset, set \textit{Production Branch} to \texttt{deploy/sprint2}, and add the variable \texttt{VITE\_API\_URL} with the value \texttt{https://}\allowbreak\texttt{<railway-domain>/api}.
This value is built into the app, so after changing it you must redeploy.
Each project has a \texttt{vercel.json} that sends every route to the main page, so reloading an inner page does not give a 404 error.

\paragraph{Last step.}
Copy the two Vercel domains into the Railway variable \texttt{CORS\_ALLOWED\_ORIGINS} (and the main one into \texttt{FRONTEND\_URL}), so the browser can call the API\@.
To switch sprints, change the branch in Railway and in both Vercel projects.

\subsection{Verification}

These checks confirm that the system works after an installation or a deployment.
The expected result is written after each check.

\paragraph{Local installation.}
Start the backend (step 2 of the installation) and run these commands in another terminal.
The first two must print \texttt{200}.
The third must print a JSON list with the demo store rooms created by the seeder.
\begin{lstlisting}[language={}]
curl -s -o /dev/null -w "%{http_code}\n" http://127.0.0.1:8000/up
curl -s -o /dev/null -w "%{http_code}\n" http://127.0.0.1:8000/api/storeRooms
curl -s http://127.0.0.1:8000/api/storeRooms
\end{lstlisting}
Then open \texttt{http://localhost:5173} (main web app) and \texttt{http://localhost:5174} (mobile web app).
Both pages must load.
The main web app must show the login page, and signing in with a seeded user must work without errors in the browser console.
An error about CORS means that the address of the web app is missing in \texttt{CORS\_ALLOWED\_ORIGINS}.
Finally, the test commands of step 5 must give the results listed there.

\paragraph{Deployed system.}
Replace the placeholders with the real domains (Table~\ref{tab:env-vars}).
\begin{enumerate}[nosep]
  \item \textbf{API.}
        \texttt{https://}\allowbreak\texttt{<railway-domain>/up} must answer \texttt{200}.
        In Railway, the last deployment must show the status \textit{Success}, and its pre-deploy step must show that the migrations ran.
  \item \textbf{Web apps.}
        Both Vercel domains must load the page.
        Reloading an inner page must not give a 404 error.
  \item \textbf{Connection.}
        Sign in from the Vercel web app with the admin user created in the first deployment.
        The browser console must show no CORS or host errors, and the store room list must load.
  \item \textbf{Uploaded files.}
        Upload a photo in a store room, redeploy the backend, and open the store room again.
        The photo must still be there.
        If it is lost, the volume at \texttt{/app/storage/app} is missing.
\end{enumerate}
